% GENERATED FILE. Edit canonical lessons, then run npm run generate:books.
% canonical-source-hash: fnv1a64:e0f685628bd4aea6
% canonical-lessons: MR-C547-anibani, MR-C547-agnisamak, MR-C547-gajar-alarm, MR-C547-suraksaraksak, MR-C547-atokyat-anne

\chapter{Emergencies}
\label{ch:mr-emergencies}
\hlchaptermodality{547}

\begin{chapteropening}
\textbf{By the end of this chapter:} \emph{I can talk about emergencies, alarms and bringing things under control.}

Complete the last lesson of chapter 547: I can talk about emergencies, alarms and bringing things under control.
\end{chapteropening}

\section[\texorpdfstring{\mr{आणीबाणी} (āṇībāṇī)}{आणीबाणी (āṇībāṇī)}]{\mr{आणीबाणी} (āṇībāṇī) — an emergency}

\label{lesson:MR-C547-anibani}



\begin{quote}
Before the new one: say the Marathi for to take an oath, then the Marathi for to dare.
\end{quote}

\subsection*{You'll want to know: \mr{आणीबाणी}}
\textbf{\mr{आणीबाणी}} — \emph{āṇībāṇī} — \textquotedblleft{}an emergency\textquotedblright{}.

\begin{quote}
\mr{आणीबाणीच्या} \mr{वेळी} \mr{हा} \mr{दरवाजा} \mr{उघडतात}.
\end{quote}

\begin{quote}
\emph{āṇībāṇīcyā veḷī hā darvājā ughaḍtāt.}
\end{quote}

\begin{quote}
This door is opened in an emergency.
\end{quote}

Feminine.

\begin{cousinweb}
\textbf{\mr{आणीबाणी}} (\emph{āṇībāṇī}) is a Marathi word for a moment of crisis; its origin is not settled.
\end{cousinweb}

\begin{grammarlens}[title={What you've built}]
One more word for saying what happens if.
\end{grammarlens}

\subsection*{Your turn}
\begin{itemize}
\raggedright
  \item \emph{Say it:} \emph{āṇībāṇī}
  \item \emph{Say it:} \emph{āṇībāṇī}, once more
  \item \emph{Say it:} \emph{āṇībāṇīcyā veḷī hā darvājā ughaḍtāt.} — This door is opened in an emergency.
  \item \emph{Recall:} say the Marathi for to take an oath, then the Marathi for to dare, then say \emph{āṇībāṇī} again
\end{itemize}

\begin{tcolorbox}[breakable,colback=teal!4,colframe=teal!35!black,title={Before you move on}]
What does \mr{आणीबाणी} mean? (\textquotedblleft{}An emergency\textquotedblright{}.) Say it once more.
\end{tcolorbox}
\section[\texorpdfstring{\mr{अग्निशामक} (agniśāmak)}{अग्निशामक (agniśāmak)}]{\mr{अग्निशामक} (agniśāmak) — a fire extinguisher}

\label{lesson:MR-C547-agnisamak}



\begin{quote}
Before the new one: say the Marathi for to dare, then the Marathi for an emergency.
\end{quote}

\subsection*{You'll want to know: \mr{अग्निशामक}}
\textbf{\mr{अग्निशामक}} — \emph{agniśāmak} — \textquotedblleft{}a fire extinguisher\textquotedblright{}.

\begin{quote}
\mr{प्रत्येक} \mr{मजल्यावर} \mr{एक} \mr{अग्निशामक} \mr{ठेवला} \mr{आहे}.
\end{quote}

\begin{quote}
\emph{pratyek majlyāvar ek agniśāmak ṭhevlā āhe.}
\end{quote}

\begin{quote}
A fire extinguisher is kept on every floor.
\end{quote}

Masculine.

\begin{cousinweb}
\textbf{\mr{अग्निशामक}} (\emph{agniśāmak}) is a tatsama coinage from Sanskrit \emph{agni}, \textquotedblleft{}fire\textquotedblright{} (Monier-Williams), and a noun of \emph{śam-}, \textquotedblleft{}to quench\textquotedblright{}.
\end{cousinweb}

\begin{grammarlens}[title={What you've built}]
One more word for saying what happens if.
\end{grammarlens}

\subsection*{Your turn}
\begin{itemize}
\raggedright
  \item \emph{Say it:} \emph{agniśāmak}
  \item \emph{Say it:} \emph{agniśāmak}, once more
  \item \emph{Say it:} \emph{pratyek majlyāvar ek agniśāmak ṭhevlā āhe.} — A fire extinguisher is kept on every floor.
  \item \emph{Recall:} say the Marathi for to dare, then the Marathi for an emergency, then say \emph{agniśāmak} again
\end{itemize}

\begin{tcolorbox}[breakable,colback=teal!4,colframe=teal!35!black,title={Before you move on}]
What does \mr{अग्निशामक} mean? (\textquotedblleft{}A fire extinguisher\textquotedblright{}.) Say it once more.
\end{tcolorbox}
\section[\texorpdfstring{\mr{गजर} (gajar)}{गजर (gajar)}]{\mr{गजर} (gajar) — an alarm}

\label{lesson:MR-C547-gajar-alarm}



\begin{quote}
Before the new one: say the Marathi for an emergency, then the Marathi for a fire extinguisher.
\end{quote}

\subsection*{You'll want to know: \mr{गजर}}
\textbf{\mr{गजर}} — \emph{gajar} — \textquotedblleft{}an alarm\textquotedblright{}.

\begin{quote}
\mr{पहाटे} \mr{पाच} \mr{वाजता} \mr{घड्याळाचा} \mr{गजर} \mr{वाजला}.
\end{quote}

\begin{quote}
\emph{pahāṭe pāc vājtā ghaḍyāḷācā gajar vājlā.}
\end{quote}

\begin{quote}
The alarm clock rang at five in the morning.
\end{quote}

Masculine.

\mr{गजर} \mr{वाजणे}, an alarm rings.

\begin{cousinweb}
\textbf{\mr{गजर}} (\emph{gajar}) is the word for an alarm or a loud call; its origin is not settled.
\end{cousinweb}

\begin{grammarlens}[title={What you've built}]
One more word for saying what happens if.
\end{grammarlens}

\subsection*{Your turn}
\begin{itemize}
\raggedright
  \item \emph{Say it:} \emph{gajar}
  \item \emph{Say it:} \emph{gajar}, once more
  \item \emph{Say it:} \emph{pahāṭe pāc vājtā ghaḍyāḷācā gajar vājlā.} — The alarm clock rang at five in the morning.
  \item \emph{Recall:} say the Marathi for an emergency, then the Marathi for a fire extinguisher, then say \emph{gajar} again
\end{itemize}

\begin{tcolorbox}[breakable,colback=teal!4,colframe=teal!35!black,title={Before you move on}]
What does \mr{गजर} mean? (\textquotedblleft{}An alarm\textquotedblright{}.) Say it once more.
\end{tcolorbox}
\section[\texorpdfstring{\mr{सुरक्षारक्षक} (surakṣārakṣak)}{सुरक्षारक्षक (surakṣārakṣak)}]{\mr{सुरक्षारक्षक} (surakṣārakṣak) — a security guard}

\label{lesson:MR-C547-suraksaraksak}



\begin{quote}
Before the new one: say the Marathi for a fire extinguisher, then the Marathi for an alarm.
\end{quote}

\subsection*{You'll want to know: \mr{सुरक्षारक्षक}}
\textbf{\mr{सुरक्षारक्षक}} — \emph{surakṣārakṣak} — \textquotedblleft{}a security guard\textquotedblright{}.

\begin{quote}
\mr{इमारतीच्या} \mr{दारात} \mr{एक} \mr{सुरक्षारक्षक} \mr{उभा} \mr{असतो}.
\end{quote}

\begin{quote}
\emph{imārtīcyā dārāt ek surakṣārakṣak ubhā asto.}
\end{quote}

\begin{quote}
A security guard stands at the door of the building.
\end{quote}

Masculine.

\begin{cousinweb}
\textbf{\mr{सुरक्षारक्षक}} (\emph{surakṣārakṣak}) is a compound built from \mr{सुरक्षा}, \textquotedblleft{}safety\textquotedblright{}, and Sanskrit \emph{rakṣaka}, \textquotedblleft{}a protector\textquotedblright{} (Monier-Williams).
\end{cousinweb}

\begin{grammarlens}[title={What you've built}]
One more word for saying what happens if.
\end{grammarlens}

\subsection*{Your turn}
\begin{itemize}
\raggedright
  \item \emph{Say it:} \emph{surakṣārakṣak}
  \item \emph{Say it:} \emph{surakṣārakṣak}, once more
  \item \emph{Say it:} \emph{imārtīcyā dārāt ek surakṣārakṣak ubhā asto.} — A security guard stands at the door of the building.
  \item \emph{Recall:} say the Marathi for a fire extinguisher, then the Marathi for an alarm, then say \emph{surakṣārakṣak} again
\end{itemize}

\begin{tcolorbox}[breakable,colback=teal!4,colframe=teal!35!black,title={Before you move on}]
What does \mr{सुरक्षारक्षक} mean? (\textquotedblleft{}A security guard\textquotedblright{}.) Say it once more.
\end{tcolorbox}
\section[\texorpdfstring{\mr{आटोक्यात} \mr{आणणे} (āṭokyāt āṇṇe)}{आटोक्यात आणणे (āṭokyāt āṇṇe)}]{\mr{आटोक्यात} \mr{आणणे} (āṭokyāt āṇṇe) — to bring under control}

\label{lesson:MR-C547-atokyat-anne}



\begin{quote}
Before the new one: say the Marathi for an alarm, then the Marathi for a security guard.
\end{quote}

\subsection*{You'll want to know: \mr{आटोक्यात} \mr{आणणे}}
\textbf{\mr{आटोक्यात} \mr{आणणे}} — \emph{āṭokyāt āṇṇe} — \textquotedblleft{}to bring under control\textquotedblright{}.

\begin{quote}
\mr{सुरक्षारक्षकाने} \mr{अग्निशामक} \mr{वापरून} \mr{आग} \mr{आटोक्यात} \mr{आणली}.
\end{quote}

\begin{quote}
\emph{surakṣārakṣakāne agniśāmak vāprūn āg āṭokyāt āṇlī.}
\end{quote}

\begin{quote}
The security guard used the extinguisher and brought the fire under control.
\end{quote}

Past, with -\mr{ने} (he …): \mr{त्याने} \mr{आग} \mr{आटोक्यात} \mr{आणली}, \emph{tyāne āg āṭokyāt āṇlī}.

\begin{cousinweb}
\textbf{\mr{आटोक्यात} \mr{आणणे}} (\emph{āṭokyāt āṇṇe}) is \mr{आटोका}, \textquotedblleft{}reach, control\textquotedblright{}, in the locative, with \mr{आणणे}, \textquotedblleft{}to bring\textquotedblright{}; the origin of \mr{आटोका} is not settled.
\end{cousinweb}

\begin{grammarlens}[title={What you've built}]
One more word for saying what happens if.
\end{grammarlens}

\subsection*{Your turn}
\begin{itemize}
\raggedright
  \item \emph{Say it:} \emph{āṭokyāt āṇṇe}
  \item \emph{Say it:} \emph{āṭokyāt āṇṇe}, once more
  \item \emph{Say it:} \emph{surakṣārakṣakāne agniśāmak vāprūn āg āṭokyāt āṇlī.} — The security guard used the extinguisher and brought the fire under control.
  \item \emph{Recall:} say the Marathi for an alarm, then the Marathi for a security guard, then say \emph{āṭokyāt āṇṇe} again
\end{itemize}

\begin{tcolorbox}[breakable,colback=teal!4,colframe=teal!35!black,title={Before you move on}]
What does \mr{आटोक्यात} \mr{आणणे} mean? (\textquotedblleft{}To bring under control\textquotedblright{}.) Say it once more.
\end{tcolorbox}
